\section{Rekursja prosta}
Klasa funkcji pierwotnie rekurencyjnych \(\prec\) to najmniejsza klasa zawierająca:
\begin{itemize}
\singlespacing
    \item \( O^k(x_1, \ldots, x_n) = 0\), zero,
    \item \( S(x) = x+1\), następnik,
    \item \( I^n_k(x_1, \ldots, x_n) = x_k\) dla \(k \leq n\), rzutowania.
\end{itemize}
i zamknięta na schemat rekursji prostej 



% zabrane z licencjat/chapters/mfi/natural_operations
\begin{theorem}[Schemat rekursji (indukcji) prostej]
    Niech dane będą zbiory \( A, Z \) oraz funkcje \( g: A \rightarrow Z, h : Z \times \natural \times A \rightarrow Z \).
    
    Wtedy istnieje dokładnie jedna funkcja \( f : \natural \times A \rightarrow Z \), dla której zachodzi
    \begin{enumerate}
        \item \( f(0, a) = g(a) \) -- warunek początkowy (bazowy)
        \item \( f(n', a) = h(f(n, a), n, a) \) -- rekursor
    \end{enumerate}
\end{theorem}

Przy pomocy tego schematu możemy zdefiniować podstawowe operacje dla \( A = Z = \natural \).

\begin{multicols}{2}
    \textbf{Dodawanie}
    \[
    \begin{cases}
        f(0, n) = g(n) = n \\
        f(k + 1, n) = h(f(k, n), k, n) = f(k, n) + 1
    \end{cases}
    \]
    Dodawanie definiujemy jako \( k + n = f(k, n) \)
\columnbreak\vfill
    \textbf{Mnożenie}
    \[
    \begin{cases}
        f(0, n) = g(n) = 0 \\
        f(k + 1, n) = h(f(k, n), k, n) = f(k, n) + n
    \end{cases}
    \]
    Mnożenie definiujemy jako \( k \cdot n = f(k, n) \)
\end{multicols}

\textbf{Odejmowanie}
\[
    x \dotminus y =
    \begin{cases}
        x - y, \text{ gdy } x \geq y \\
        0, \quad\quad \text{ gdy } x < y
    \end{cases}
\]
Definiujemy funkcję pomocniczą \(g(x) = x \dotminus 1\):
\begin{align*}
    g(0) & = 0 & g(x+1) & = x
\end{align*}
Teraz możemy zdefiniować \(x \dotminus y\) jako:
\begin{align*}
    x \dotminus 0 & = x & x \dotminus (y+1)& = x \dotminus y \dotminus 1
\end{align*}